\chapter{State of the Art}\label{ch:related-work}

% DRAFT NOTES FOR THE AUTHOR (comments only; they do not appear in the PDF)
% - Every statement about another tool was checked on 2026-11-17 against the tool's own paper, preprint or README; where only the abstract (or a
%   third-party description) could be read, the sentence carries a "% TODO verify" comment. The text paraphrases; it contains no direct quotations.
% - Numbers are kept to a minimum. The quantitative comparison with these tools is Chapter~\ref{ch:evaluation} (HG2K_RESULTS.md, section 8).
% - None of the tools described here is a component of the tool of this thesis; they are related work and baselines.
% - Section numbers cited in comments below refer to the arXiv versions that were read, which can differ from the published layout.

This chapter describes what has been built to keep software runnable when its environment breaks, and where the existing work stops. It has no results of its own. Section~\ref{sec:sota-repro} recalls why the problem matters. Sections~\ref{sec:sota-classical} and~\ref{sec:sota-llm} describe two families of tools, the older rule-based and \ac{KG}-based ones and the newer ones that use a \ac{LLM}. Section~\ref{sec:sota-bench} names the benchmarks on which the tools are measured. Section~\ref{sec:sota-gap} states the gap that the thesis addresses. The numbers with which this thesis compares itself to the tools are given in Chapter~\ref{ch:evaluation}.

\section{The Reproducibility Problem}\label{sec:sota-repro}

The motivation for automatic repair comes from large measurement studies. Pimentel et al.\ executed notebooks mined from GitHub; of 863{,}878 attempted executions of valid notebooks, only 24.11\,\% ran without errors and only 4.03\,\% produced the same results, and missing dependencies were among the most common causes of failure \cite{pimentel2019large}. Samuel and Mietchen repeated the exercise for notebooks that accompany biomedical publications, the corpus on which this thesis is evaluated: of 27{,}271 notebooks, 10{,}388 had all their declared dependencies installed and were re-run, and 1{,}203 of these ran without errors \cite{samuel2024computational}. They attribute the failures mainly to dependencies. For machine-learning notebooks, Jin et al.\ report that rebuilding the original environment by downgrading libraries to the submission date does not restore reproducibility \cite{jin2026modernization}.

These studies measure the decay of environments; they do not repair it. The tools below are the attempts to do so automatically. They differ in what they take as input (a single file, a notebook or a whole repository), in what they produce, and in whether the result is checked by running it.

\section{Rule-Based and Knowledge-Graph Approaches}\label{sec:sota-classical}

The first family infers a configuration from the program text and a pre-built knowledge base about packages. These tools do not use a \ac{LLM}. Table~\ref{tab:sota-tools} summarises them together with the tools of Section~\ref{sec:sota-llm}.

\subsection{Import scanning: pipreqs}

The simplest tool is \texttt{pipreqs}, an open-source command-line program that scans the imports of a project and writes a \texttt{requirements.txt} file \cite{pipreqs}. According to its documentation it can pin exact versions, allow compatible ones or leave them open, and by default it queries the package index unless told to use only local information. Its documentation describes no reasoning about the Python interpreter or about conflicts between the listed packages, and the related-work discussion of PyEGo \cite{ye2022pyego} likewise counts it as a tool that handles third-party packages only. It serves in the literature as the naive baseline; no academic paper about it is known to us.% TODO verify: that no academic paper about pipreqs exists (not searched for).

\subsection{DockerizeMe}

DockerizeMe by Horton and Parnin targets code snippets \cite{horton2019dockerizeme}. It builds a knowledge base offline for the 10{,}000 highest-ranked packages of the \ac{PyPI}, from static analysis of the packages, from installing them and parsing the resulting error messages, and from association rules mined from public Dockerfiles, which supply the operating-system packages. The base is stored as a graph. At inference time DockerizeMe parses the snippet for its imports, discards the standard library, maps the remaining import names to packages and adds the transitive dependencies in an install order. The output is a Dockerfile. The inference yields a single static configuration; running the snippet is part of the evaluation, not of the inference. The tool does not infer the Python version: the evaluation uses Python 2.7, and snippets that need another version are stated to be unsupported. The same authors had earlier built Gistable, a data set of Python gists from GitHub \cite{horton2018gistable}, and DockerizeMe defines the hard subset of it, HG2.9K, on which most later tools are measured (Section~\ref{sec:sota-bench}).% Sections read (arXiv:1905.11127v1): III, IV, V, VI, VII-A, VII-D, VIII. Python 2.7.14 image in VII-A.

\subsection{PyEGo}

PyEGo by Ye et al.\ extends the inference to the Python interpreter and to system libraries \cite{ye2022pyego}. The authors first show empirically that many programs depend on a specific interpreter version or on system libraries, and that earlier tools, which concentrate on packages, miss both. PyEGo builds a knowledge graph, PyKG, of the relations and constraints among packages, interpreter versions and system libraries. For a program it extracts features, queries the graph for candidates of the three kinds, forms a program-specific sub-graph and solves the version constraints in it; the result is the latest compatible set of versions. PyEGo is evaluated on HG2.9K, on 100 larger Python projects and on 4{,}836 notebooks \cite{ye2022pyego}. In the evaluation of its authors, a program counts as handled when the import errors disappear.

\subsection{PyCRE and ReadPyE}

Cheng et al.\ proposed two successive knowledge-graph tools. PyCRE builds an ontology and graphs for more than 10{,}000 popular packages, separately for Python~2 and Python~3, matches the modules that a file uses to candidate libraries by a matching degree over their reference paths, and resolves version conflicts with a heuristic graph traversal that falls back to a satisfiability solver \cite{cheng2022pycre}. Its input is a single file that must parse; its output is a \texttt{requirements.txt} and a Dockerfile that records the inferred Python version. ReadPyE, the later tool, builds an ecosystem knowledge graph automatically, matches candidate packages to unknown modules with a naming-similarity measure, treats the choice of candidates as an optimisation problem and, according to its abstract, validates the inferred environment iteratively and adjusts it using exception templates found in the validation logs \cite{cheng2024revisiting}.% TODO verify: the description of ReadPyE rests on the abstract (a third-party copy) and on the description in PLLM (arXiv:2501.16191, Section II-A); the journal paper was not accessible. Input, output and evaluation data sets of ReadPyE are therefore not described.
Whenever this thesis quotes a ReadPyE success rate on HG2.9K, it is the value measured by the PLLM authors, not the value reported by the ReadPyE authors.

\subsection{SnifferDog and later work}

SnifferDog addresses Jupyter notebooks specifically \cite{wang2021sniffer}. It builds a database of the interfaces of Python libraries in many versions, analyses the imports and the used interfaces of a notebook to find candidate libraries and versions, then installs the candidates and checks whether the notebook executes and, ideally, reproduces its stored results. Its authors report that its database, built from about 488 libraries, did not cover the libraries of many notebooks in their ground-truth set, which they had to exclude; and that remaining failures stem partly from Python-version compatibility. PyEGo's authors note that SnifferDog does not handle system libraries or the interpreter \cite{ye2022pyego}.% Verified in the SnifferDog PDF (arXiv:2103.02959): API bank, installation and execution check, "385 of 725 notebooks" outside the bank, installation failures linked to the selected Python version.

A recent member of the family is PCREQ, which infers compatible requirements after a library has been upgraded, by combining a version analysis with a code-compatibility analysis, and which, according to its abstract, outputs a requirements file and a repair report \cite{lei2025pcreq}. It is evaluated on a different benchmark of upgrade scenarios, not on HG2.9K, so its numbers are not comparable with those of the other tools.% TODO verify: the venue of PCREQ (the abstract page says ACM TOSEM, the HTML header says JACM) and what the "repair report" contains; only the abstract and the introduction were read.

\subsection{What the family has in common}

These tools need no model at inference time and, in the measurements of the PLLM authors, were faster per file than PLLM; they work best where the problem is a lookup, such as an import name that differs from its package name. Their successors name the shared weaknesses: the knowledge base must be kept up to date, and it covers the packages that its authors processed (SnifferDog's own evaluation excluded notebooks whose libraries were outside its database). They struggle with complex projects, according to the authors of PLLM \cite{bartlett2025lastdependency}. The outputs described in the papers are descriptions of an environment, namely a requirements file or a Dockerfile, and none of the papers describes editing the source code of the program. Whether the result is verified by running varies: DockerizeMe infers one configuration without feedback, SnifferDog installs and executes its candidates, and ReadPyE iterates on the validation logs. In none of these papers is the result reported together with the evidence on which each individual choice rests.% TODO: author to confirm -- the last sentence is a statement about what the papers do not contain, based on reading the cited papers for their output formats; not every section of each paper was read.

\section{LLM-Based and Agentic Approaches}\label{sec:sota-llm}

\subsection{PLLM: inference with a language model}

PLLM by Bartlett et al.\ replaces the knowledge graph by a \ac{LLM} \cite{bartlett2025lastdependency}. Given one Python file, the model proposes the Python version and the modules with their versions; with \ac{RAG}, the versions that exist for each module are retrieved from the package index, and the system builds a Dockerfile from the proposal. If the build or the run fails, the error is classified and a prompt specific to the error type asks the model for a repair, in a loop with a fixed limit. The paper selects the model Gemma2 with nine billion parameters and \ac{RAG}, and counts a file as fixed when it runs without import, attribute and syntax errors. PLLM is evaluated on HG2.9K and compared with PyEGo and ReadPyE. Its authors observe that the two knowledge-graph tools still fix a subset of the files that PLLM does not, describe the two approaches as complementary and name hybrid architectures that combine prompt-based inference with knowledge graphs as future work (Section~VII of the arXiv version). This observation is a motivation for this thesis (Chapter~\ref{ch:introduction}). The paper addresses single files; in the version read, it does not discuss explaining the fixes to a user, recording their provenance or reporting a confidence.% Verified in arXiv:2501.16191v2: Sections III (method), IV-C/IV-D (loop value, success definition), V-A (Table III), VII (conclusion and future work). Section VI (threats to validity) is deliberately not used. Venue: none stated in the paper; cited as an arXiv preprint. TODO verify: venue (a workshop acceptance is mentioned in the vision document but not in the paper).

\subsection{Agents for whole repositories}

A second line of work lets an agent act on a whole repository. Repo2Run is presented by its authors as the first \ac{LLM}-based agent that builds an executable environment for arbitrary Python repositories \cite{hu2025repo2run}. It works in a container, keeps snapshots so that a failed command can be rolled back, and synthesises a Dockerfile from the commands that finally succeeded. Success means that the repository's unit tests can be executed, whether or not they pass, and the benchmark contains only repositories that come with unit tests: 420 of them, on which the authors report a success rate of 86.0\,\%. EnvBench is a benchmark of 329 Python and 665 repositories in languages of the Java Virtual Machine, chosen because they pose genuine configuration problems; its best approach, an agent with a single shell tool and GPT-4o, configures 6.69\,\% of the Python repositories \cite{eliseeva2025envbench}.% Verified in arXiv:2502.13681v4 (abstract, Sections 3, 4.1-4.4) and arXiv:2503.14443v1 (abstract, Sections 3.1, 3.3, 5). TODO verify: the venue of both (neither arXiv version read states one; the PDF header of EnvBench is reported in the bib comment as ICLR 2025, Repo2Run as NeurIPS 2025).
The two figures are not comparable with each other or with HG2.9K: Repo2Run counts a repository as built when its tests can run, whereas EnvBench's Python measure is a static check, with a type checker, for missing imports after the setup attempt; the tests are not run \cite{hu2025repo2run,eliseeva2025envbench}. What the pair shows is that environment construction for real repositories remains hard, in particular for repositories chosen for their difficulty, and that the field optimises whether an environment works.

\subsection{Knowledge graphs and verification in neighbouring tasks}

Two recent papers combine a model with structured knowledge or with a verifying tool for tasks that are not dependency repair, and are relevant to the design in this thesis only as analogies. KGCompass builds a knowledge graph that links issues and pull requests to files, classes and functions and gives the paths in it to a \ac{LLM} to localise and fix bugs in repositories; it is evaluated on software issues, not on environments \cite{yang2025kgcompass}. TerraRepair repairs security misconfigurations in Terraform infrastructure code with an agent that retrieves the dependency context, checks the installed provider schema and re-runs the scanner before it returns a candidate; its authors identify fixes that a model claims but that do not hold as a key failure mode of unaided model repair \cite{mengistu2026terrarepair}. Their domains are different from ours; the point taken from them is that grounding a model's proposal in retrieved facts and checking it with a tool is an established pattern.% KGCompass (arXiv:2503.21710): software issue repair on SWE-bench Lite -- NOT dependency/environment repair; the task brief described it as the repo-scale incarnation of the PyEGo/ReadPyE idea, which the paper does not support. TODO verify: venue (arXiv only on both pages read). TerraRepair (arXiv:2607.11390): Terraform; the arXiv comments field says accepted at ESEM 2026, TODO verify.

\subsection{Repairing the code, not only the environment}

Dependency problems are not always solved by choosing other versions: a newer library may have removed an interface that the code uses. For notebooks, Zhu et al.\ propose to restore executability by automatically upgrading deprecated interfaces in the code \cite{zhu2021relancer}, and Jin et al.\ modernise machine-learning notebooks for reproducibility \cite{jin2026modernization}. Both treat the code as part of what is repaired.% TODO verify: only the titles and bibliographic data of both papers were confirmed; the first paper's body could not be read (IEEE page empty), so its method, tool name, data set and limitations are not described. Jin et al. are cited here only for the title-level claim and for the finding already used in Chapter 1.
The tools of Section~\ref{sec:sota-classical} infer an environment for the code as it is, and the repair cycle of PLLM infers missing modules, different versions or additional dependencies \cite{bartlett2025lastdependency}.% Verified in arXiv:2501.16191v2, Section III-E. TODO verify: whether PLLM ever edits the source code was not checked beyond this description.

\begin{table}[ht]
  \centering
  \caption{The tools of Sections~\ref{sec:sota-classical} and~\ref{sec:sota-llm}: input, output and basis of the decision, as stated in the cited papers. Success rates are not given here; see Chapter~\ref{ch:evaluation}.}
  \label{tab:sota-tools}
  \small
  \begin{tabular}{@{}>{\raggedright\arraybackslash}p{0.17\linewidth}>{\raggedright\arraybackslash}p{0.16\linewidth}>{\raggedright\arraybackslash}p{0.25\linewidth}>{\raggedright\arraybackslash}p{0.34\linewidth}@{}}
    \hline
    Tool & Input & Output & Basis of the decision \\
    \hline
    pipreqs \cite{pipreqs} & project imports & \texttt{requirements.txt} & scan of imports; package index lookup \\
    DockerizeMe \cite{horton2019dockerizeme} & single snippet & Dockerfile (Python 2.7) & offline knowledge base of 10{,}000 packages; import-to-package lookup; transitive dependencies \\
    PyEGo \cite{ye2022pyego} & single file, project or notebook & compatible versions of packages, interpreter and system libraries & knowledge graph PyKG; constraint solving on a program-specific sub-graph \\
    PyCRE \cite{cheng2022pycre} & single file that parses & \texttt{requirements.txt} and Dockerfile with a Python version & ontology and knowledge graphs for Python~2 and~3; matching of reference paths; graph traversal with a solver fallback \\
    ReadPyE \cite{cheng2024revisiting} & not stated in the abstract & not stated in the abstract & ecosystem knowledge graph; naming similarity; optimisation; validation against exception templates \\
    SnifferDog \cite{wang2021sniffer} & notebook & candidate libraries and versions, checked by execution & database of library interfaces by version; analysis of imports and interface use \\
    PLLM \cite{bartlett2025lastdependency} & single file & Dockerfile with Python version and modules & \ac{LLM} with retrieval of package versions; loop on error messages \\
    Repo2Run \cite{hu2025repo2run} & repository with unit tests & Dockerfile & \ac{LLM} agent in a container with rollback; the repository's tests must be executable \\
    \hline
  \end{tabular}
\end{table}

\section{Benchmarks and Data Sets}\label{sec:sota-bench}

\paragraph{Gistable and HG2.9K.}
Gistable is a collection of 10{,}259 Python gists from GitHub in which, in a clean Python~2.7 environment, only 24.4\,\% ran and 52.4\,\% failed with an import error \cite{horton2018gistable}. The DockerizeMe paper defines HG2.9K as the 2{,}891 gists that still fail with an import error after Gistable's naive inference of packages from import names \cite{horton2019dockerizeme}. The set is used by PyEGo, ReadPyE and PLLM and is the secondary benchmark of this thesis; its gists are single files, so it tests the core of dependency repair at scale but not repository-level or notebook-specific problems. Because it was built when Python~2 was still common, many of its gists are written for Python~2; Chapter~\ref{ch:evaluation} measures their share.% TODO verify: the venue and pages of the Gistable paper (taken from the reference list of the DockerizeMe paper; the arXiv page shows no journal reference).

\paragraph{The GigaScience corpus.}
The notebooks of Samuel and Mietchen, described in Section~\ref{sec:sota-repro}, are the primary data set of this thesis \cite{samuel2024computational}. They come from real publications, carry no tests and fail for many reasons that are not dependency problems.

\paragraph{Other resources.}
The notebooks of Pimentel et al.\ form a large corpus of GitHub notebooks (1{,}159{,}166 unique notebooks from 264{,}023 repositories in the study) \cite{pimentel2019large}; EnvBench and the benchmark of Repo2Run target repositories \cite{eliseeva2025envbench,hu2025repo2run}. They are not used in this thesis.

\section{Positioning: The Gap}\label{sec:sota-gap}

Taken together, the reviewed work leaves four properties that, to the knowledge of the author, no single tool offers at once.

\begin{enumerate}
  \item \textbf{Repository scale.} The knowledge-graph tools and PLLM work on single files, and SnifferDog on single notebooks. Repo2Run and the EnvBench approaches work on repositories, but they assume or measure something else (runnable tests, a static import check).
  \item \textbf{Interpreter, packages and code.} The classical tools infer an environment for code that stays as it is. The work on repairing the code of notebooks that was reviewed concerns deprecated interfaces and does not choose the environment.% TODO verify: for Zhu et al. this rests on the title only (paper not readable).
  \item \textbf{Verification by execution.} Some tools check their result by running it (SnifferDog, ReadPyE, PLLM and Repo2Run in different forms). In a repair loop that follows the principle ``the model proposes, the check decides'', each change is accepted only after the target has actually run.
  \item \textbf{Transparency.} None of the papers reviewed presents, for each repair step, whether it rests on package-index metadata or on a model's suggestion, whether it was verified, and how much confidence it deserves. The tools report outcomes; the reasoning is discarded. PCREQ emits a repair report, whose content was not examined.% TODO verify: content of the PCREQ repair report; if it records per-change evidence, soften this paragraph.
\end{enumerate}

The thesis addresses this combination for failures of Python notebooks and repositories, and states its contribution accordingly: the contribution is the transparency of the repair and the measurement of how much of the problem is within reach of dependency repair, not a record success rate. On HG2.9K the tool is mid-table among the tools, and that is expected; Chapter~\ref{ch:evaluation} reports the comparison with its caveats. The gap corresponds to the research questions of Chapter~\ref{ch:introduction}: the first (repair) asks how well an agentic, model-based approach repairs notebooks, the second (scope) asks what share of failures dependency repair can address at all, and the third (transparency) asks how the process can be made transparent. The fourth question, whether explanations improve trust, is future work.
